\section{Introduction}
\label{sec:intro}

A language model used over a long period has to learn new tasks without losing old ones. Low-rank
adapters \citep{hu2022lora} remove half of the problem: the pretrained weights stay frozen, each task
adds a small update, and an old adapter can be kept unchanged. What remains is which adapters should
act on a given input. If the model is told the task, nothing is forgotten, but in practice it is not
told. If all adapters are left on, each new adapter changes what the model computes on earlier inputs
(Section~\ref{sec:results}: with five tasks, plain adapters keep \FivePlainAP{} of
\FivePlainLearned{} Rouge-L points).

Published methods restrict what a new adapter may learn, through orthogonal or interference-free
subspaces \citep{wang2023olora,liang2024inflora,spacelora2026}, or add modules that weight the
adapters per input: learned gates or attention over adapters
\citep{liang2025gainlora,zhao2024sapt}, or a mixture model over frozen embeddings
\citep{latentlora2026}. They reduce forgetting and are evaluated on benchmarks whose inputs begin with
the instruction of the task. To our knowledge none comes with a statement that the answer to a given
input \emph{will not} change when the next task is learned: a restriction is approximate, a learned
gate has to be protected itself, and a router that picks the best match can re-decide an old input
when a new task appears (on machine paraphrases of real tasks this proved rare: Section~\ref{sec:paraphrase}).

The idea of this paper is to give every task one frozen set of adapters and a switch that sends each
input to one task; we call the system FK-LoRA, after the rule \emph{finders keepers}. An answer can
then change only if its route changes (Lemma~\ref{lem:route}), so retention is a property of the rule
that sets the route, and it can be proved or measured for that rule. We give three rules
(Section~\ref{sec:method}).

\paragraph{Without instructions: \recog{} and \handicap{}.} The instruction gives the task away.
Without it, every method loses at least the conflict between two tasks, the inputs that occur in both
with different right answers (Proposition~\ref{prop:noinstructionprice}); the conflict is zero when no
input occurs in two tasks, as in our benchmarks. Without instructions, the two published methods we ran
changed \AudSeedsGainloraChanged{} (GainLoRA) and \AudSeedsInfloraChanged{} (InfLoRA) of
\AudSeedsCompared{} answers, against \FiveGainChanged{} and \FiveInfChanged{} of \FiveGainCompared{}
with them (Section~\ref{sec:nodefpublished}). \recog{} finds the task of an input from the frozen
model's own states, with one vector per task and one matrix for all tasks; it is the same for every
order of arrival and blocks no task. Under \handicap{} an older task keeps what it found. An audit
estimates from held-out inputs the share of changed answers (Proposition~\ref{prop:audit}), and the
handicap bounds in advance what later tasks can take (Theorem~\ref{thm:handicap}). \recog{} gives all
but \AudRouteFiveRidgeNotOwn{} of the \AudRouteFiveRidgeN{} test inputs of the five tasks to their own
task; on the 15 benchmark tasks, which were development data, it gives \AudDevLLamOneEnd{}\% of the test inputs to their own task, \AudScaleLongThreeEnd{}\% on the 15 tasks of a second benchmark (read once, against a line fixed
in advance), \AudFreshFourthLEnd{}\% on 30 unseen tasks, where it passed the line fixed in advance, and \AudScaleAllHundredTwentyEnd{}\% with 120 tasks (routes only;
Section~\ref{sec:recogniser}). With trained answers, plain adapters behind it change
\AudSeedsRidgeChanged{} of the \AudSeedsCompared{} answers, and beat the published methods in this
setting by a line fixed before the run (Section~\ref{sec:nodeftrained}). The caveats belong to the
result: with a tenth of their epochs most of the gap in the final score is learning, and the five tasks are among the 15 development tasks on which the recogniser was chosen; with the
published number of epochs (one seed) the two methods learn the tasks and then lose them, and the gap
is forgetting. With their definitions left in, they forget far less: GainLoRA \AudFullDefOriginalGainloraFTOne{} points, with
\AudFullDefOriginalGainloraChanged{} of \AudFullDefCompared{} answers changed, and InfLoRA \AudFullDefOriginalInfloraFTOne{} points, with \AudFullDefOriginalInfloraChanged{}. In this setting the
instruction is what held most of their answers.

\paragraph{With instructions: \cert{}.} Each task claims a protected zone of input space when it is
learned, computed from the frozen word-vectors of its instruction; nothing is trained for it and it is
never changed. An input inside a zone belongs to the oldest task that claims it and is answered with
the adapters of that task. Because later tasks can only add newer claims, that answer cannot change,
and the system returns a certificate with it. Inputs outside every zone go to the nearest zone, without
a certificate. On SuperNI with T5-base no answer to an earlier task changed (\FiveCertChanged{} of
\FiveCertCompared{} on five tasks, three seeds; \FifteenCertChanged{} of \FifteenCertCompared{} on
fifteen). The price is paid in the open. A new task whose inputs fall inside an older zone cannot be
learned on those inputs, and this is counted when the task arrives; in our benchmark sequences no task collided, but in a collection of \ManyTasks{}
real tasks about one in six collided (Section~\ref{sec:many}). A certified answer can never improve
through later tasks. The guarantee covers only the inputs inside a zone: all inputs with the original
instruction, all lightly reworded inputs of the benchmark tasks, about a fifth of the medium and heavy ones, and under machine paraphrases at scale \WidenRefTfivePLow{}\% to \WidenRefTfivePHigh{}\%
(Section~\ref{sec:paraphrase}); five attempts to widen the zones missed their bars.

\paragraph{The two rules together, and a third.} The zones read the instruction, and a reworded
instruction can leave them; \recog{} reads only the example, so its routes are the same for every
wording. With instructions on the five tasks (three seeds), plain adapters behind \recog{} end at
\AudUThreeHeavyRidgeAPOne{} under a heavy rewording where the zones end at
\AudUThreeHeavyCertAPOne{}, and change at most \AudUThreeMaxChangedPerSeed{} of \AudUCompared{} answers
per seed where the zones change none. With the original instructions the zones end higher
(\AudUThreeOriginalCertAPOne{} against \AudUThreeOriginalRidgeAPOne{}) and change no answer where the
recogniser changes \AudUThreeOriginalRidgeChanged{} of \AudUThreeCompared{}
(Section~\ref{sec:exampleroutes}). The third rule builds the zones from the inputs of a task and not
from its instruction; its bound (Theorem~\ref{thm:bound}) does not grow with the number of later tasks and held,
but in our first tests frozen word-vector features made it block most later tasks (Section~\ref{sec:noinstruction}).

\paragraph{Contributions.}
\begin{itemize}
  \item \textbf{Principle and price.} With one frozen set of adapters per task, retention is a
  property of the rule that sets the route (Lemma~\ref{lem:route}); without instructions every method
  loses at least the conflict of two tasks (Proposition~\ref{prop:noinstructionprice}).
  \item \textbf{A rule without instructions that blocks no task}, with an audit and a handicap
  (Proposition~\ref{prop:audit}, Theorem~\ref{thm:handicap}); the handicap bound is loose. Two tasks built on the same texts split their inputs.
  \item \textbf{The published methods without their instructions}: both hold their answers less well than with them, and plain adapters behind \recog{} or \handicap{} beat them in this setting, both verdicts by lines fixed before the runs; compared at a tenth of
  and at the full number of epochs (Sections~\ref{sec:nodefpublished} and~\ref{sec:nodeftrained}).
  \item \textbf{A switch with a proof} for inputs that carry an instruction
  (Theorem~\ref{thm:writeonce}, Proposition~\ref{prop:edits}); on the benchmarks simpler rules that recognise the instruction reach the same scores, so the switch adds the proof and the certificate, not points; its price measured at scale
  (Sections~\ref{sec:five}, \ref{sec:fifteen}, \ref{sec:many} and~\ref{sec:paraphrase}).
  \item \textbf{An evaluation that can show the difference:} answers that change are counted, every
  method is tested with reworded instructions or without them, learning and forgetting are reported
  separately, and the line for a verdict is fixed before the run (Section~\ref{sec:protocol}).
  \item \textbf{Further tests of the zones:} a constructed look-alike case (Section~\ref{sec:stress}),
  mechanical edits (Section~\ref{sec:edits}), orders of arrival (Section~\ref{sec:orders}), unrelated
  inputs (Section~\ref{sec:background}), eight embedding tables and small stand-alone models
  (Section~\ref{sec:models}), a variant for sibling tasks (Section~\ref{sec:exactfirst}), and zones built from the inputs of a task (Theorem~\ref{thm:bound}, Section~\ref{sec:noinstruction}).
\end{itemize}

\paragraph{What we do not claim.} On these benchmarks the unedited instruction identifies the task: a
rule without training recognises the task of every benchmark test input and of \ManyNearRightPct{}\% of
the inputs of \ManyTasks{} real tasks, so ``no forgetting'' follows from recognising the instruction.
A test of retention has to take the instruction away or reword it; we do the first for two published
methods in a reduced setting (T5-base instead of the T5-large of the GainLoRA paper, whose language
setting we use; the InfLoRA paper itself reports image tasks; three seeds at a tenth of the published
epochs, and one seed at the published number, without and with instructions)
and the second at scale by machine. The first is a first look; it cannot show what other methods do, and we do not claim that
published methods in general only look good because of the instruction; rewordings written by people remain the open experiment. We do not claim that forgetting
is solved for every input: an input outside all zones has no guarantee, and under \recog{} the share
of changed answers is estimated and bounded, not ruled out. We do not claim that the gap at a tenth of
the epochs is retention (with the published number it is forgetting, in one seed), nor a result in the published setting: all GPU experiments use a smaller model, all but a few a shorter training budget, and the 15-task comparison one seed. We do not claim the recogniser as new: it is a known kind of classifier on frozen features, and what is new is what is said and measured about its routes. Recognising a task from frozen word-vectors without a trained router is not new either \citep{latentlora2026}; for the zones we add
the rule that a claim cannot be taken away, its proof and the certificate. We do not claim that the
certified switch routes reworded inputs better than simpler rules: at scale a plain nearest-signature
rule gave more of them to the right task, because the fallback of our rule was chosen on the
benchmark; a variant found at the end closes most of the gap, with the same benchmark routes and so the same GPU scores, but the best rule without zones is still
\RerouteZonesCost{} to \RerouteZonesCostEs{} points better, depending on the pivot language. We do not claim that the guarantee is free: in a large collection it refuses tasks that resemble older ones. The rewordings of the five tasks
were written by machine and looked at during development; the paraphrases at scale (back-translation
through German, French, Spanish and Italian of \WidenTasksEval{} real tasks, routes only) are not
rewordings by people. The closest published method, Latent-LoRA, has no public code, and our
re-implementation learned only the first of the five tasks, so we compare only its router
(Appendix~\ref{app:latent}).
